%!TEX root = ../QuantumNoiseAndDecoherence.tex
% ──────────────────────────────────────────────────────────────
%  Lecture 7 — Lindblad Equations I: From Microscopic Dynamics
%               to the Born–Markov Master Equation
% ──────────────────────────────────────────────────────────────

\section{Lindblad equations}\label{sec:lindblad}

Over the past lectures we assembled all the ingredients needed to
describe a system~$S$ interacting with an environment~$E$: completely
positive maps, Fock space, Gaussian models.  Our goal now is to derive
an \emph{equation of motion} for the system alone---one that does not
require tracking every degree of freedom of the environment.

% ──────────────────────────────────────────────────────────────
\subsection{Setup: system plus environment}%
\label{sec:sys-env-setup}

The total Hilbert space is
$\hilb = \hilb_S \tp \hilb_E$,
where $\hilb_S$ is typically finite-dimensional and $\hilb_E$ models
the environment---often an infinite collection of harmonic oscillators.
The joint Hamiltonian is
\begin{equation}\label{eq:H-total}
  \eqbox{H = H_S \tp \One_E + \One_S \tp H_E + V}\,,
\end{equation}
with free part $H_0 = H_S \tp \One_E + \One_S \tp H_E$ and
interaction~$V$.  Any Hamiltonian on $\hilb_S \tp \hilb_E$ admits
such a decomposition.

\begin{intuition}[Why the full description fails]
  The exact reduced state
  $\rho_S(t) = \tr_E\!\bigl(U(t)\,(\rho_S(0)\tp\rho_E)\,U(t)^\adj\bigr)$
  requires the joint unitary $U(t)=e^{-iHt}$, which couples the few
  degrees of freedom of~$S$ to the (infinitely many) degrees of
  freedom of~$E$.  We need a description involving only~$S$.
\end{intuition}

% ──────────────────────────────────────────────────────────────
\subsection{Five assumptions toward a master equation}%
\label{sec:assumptions}

\subsubsection{Assumption 1: initially decoupled state}

At $t=0$ the system and environment are in a product state:
\begin{equation}\label{eq:product-init}
  \eqbox{\rho(0) = \rho_S(0) \tp \rho_E}\,,
\end{equation}
justified when $S$ has been insulated from~$E$ by design or the
coupling is weak.  If relaxed, the reduced dynamics need not even be
completely positive.

\subsubsection{Assumption 2: weak coupling}

$V$ is a \emph{small} perturbation of~$H_0$, controlled by a
dimensionless coupling constant (e.g.\ the fine-structure constant
$\alpha\approx 1/137$).

\begin{remark}
  One may always decompose the interaction-picture coupling as
  $\tilde{V}(t) = \tilde{V}_S(t)\tp\One_E + \tilde{V}_{SE}(t)$
  with $\tr_E(\tilde{V}_{SE}(t)\,\rho_E)=0$, by absorbing any
  system-only part into a redefined~$H_S$.  This ensures the
  first-order contribution vanishes after tracing over~$E$.
\end{remark}

% ──────────────────────────────────────────────────────────────
\subsection{The interaction picture}%
\label{sec:interaction-picture}

Define $U_0(t)=e^{-iH_0 t}$ and set
\begin{equation}\label{eq:int-picture}
  \tilde{V}(t) = U_0(t)^\adj V\, U_0(t)\,,\qquad
  \tilde{\rho}(t) = U_0(t)^\adj\rho(t)\, U_0(t)\,.
\end{equation}
The equation of motion in the rotating frame involves only the
interaction:
\begin{equation}\label{eq:von-neumann-int}
  \eqbox{\frac{\dd}{\dd t}\tilde{\rho}(t)
  = -i\bigl[\tilde{V}(t),\,\tilde{\rho}(t)\bigr]}\,.
  \tag{$\star$}
\end{equation}

Integrate from~$0$ to~$t$ and substitute back into~$(\star)$ to
obtain an equation correct to second order in~$V$:
\begin{equation}\label{eq:second-order}
  \frac{\dd}{\dd t}\tilde{\rho}(t)
  = -i\bigl[\tilde{V}(t),\,\tilde{\rho}(0)\bigr]
  - \int_0^{t}\!\dd t_1\;
    \bigl[\tilde{V}(t),\,
          \bigl[\tilde{V}(t_1),\,\tilde{\rho}(t_1)\bigr]\bigr]\,.
\end{equation}
Iterating further produces the \textbf{Dyson series}; we stop at
second order because the first-order term will vanish after tracing
over~$E$ and all leading physical effects (dissipation, decoherence)
appear at this order.

% ──────────────────────────────────────────────────────────────
\subsection{Tracing out the environment}%
\label{sec:trace-out-env}

Taking $\tr_E$ of~\eqref{eq:second-order} and using
$\tr_E(\tilde{V}_{SE}(t)\,\rho_E)=0$ to kill the first-order term:
\begin{equation}\label{eq:pre-born}
  \frac{\dd}{\dd t}\tilde{\rho}_S(t)
  = -\int_0^{t}\!\dd t_1\;
    \tr_E\!\Bigl[\tilde{V}(t),\,
      \bigl[\tilde{V}(t_1),\,\tilde{\rho}(t_1)\bigr]\Bigr]\,.
\end{equation}
This is exact to second order but still involves the full
system--environment state inside the integral.

% ──────────────────────────────────────────────────────────────
\subsection{The Born approximation (Assumption 3)}%
\label{sec:born-approx}

Replace the full state inside the integral by a product:
\begin{equation}\label{eq:born}
  \eqbox{\tilde{\rho}(t_1)
  \;\approx\;
  \tilde{\rho}_S(t_1) \tp \rho_E}\,.
\end{equation}

\begin{intuition}[Why neglecting entanglement is justified]
  The Born approximation looks drastic: it discards all
  system--environment entanglement.  But we \emph{never interrogate
  the environment}.  A photon emitted by the system flies away at the
  speed of light and never returns; since we cannot perform a Bell
  test between $S$ and~$E$, the correlations are invisible to any
  system-only measurement.  The approximation breaks down if the
  environment can ``talk back'' (e.g.\ photons bouncing between
  mirrors).
\end{intuition}

Substituting into~\eqref{eq:pre-born}:
\begin{equation}\label{eq:born-eq}
  \frac{\dd}{\dd t}\tilde{\rho}_S(t)
  = -\int_0^{t}\!\dd t_1\;
    \tr_E\!\Bigl[\tilde{V}(t),\,
      \bigl[\tilde{V}(t_1),\,
            \tilde{\rho}_S(t_1)\tp\rho_E\bigr]\Bigr]\,.
\end{equation}

% ──────────────────────────────────────────────────────────────
\subsection{The Markov approximation (Assumptions 4 and 5)}%
\label{sec:markov-approx}

\textbf{Assumption 4} (dense bath spectrum).  The environment has
many closely spaced energy levels in every interval $\Delta E$.

\begin{intuition}[Continuous spectra and irreversibility]
  Only baths with absolutely continuous spectrum exhibit irreversible
  information loss---a consequence of the \emph{RAGE theorem}.
  Discrete spectra lead to Poincar\'e recurrences.
\end{intuition}

Environmental correlations therefore decay on a time scale $\tau_E$
much shorter than the system relaxation time, permitting the
replacement $\tilde{\rho}_S(t_1)\approx\tilde{\rho}_S(t)$ inside
the integral.

\textbf{Assumption 5} (Markov).  The bath spectrum is truly
continuous, so we may extend the integration:
\begin{equation}\label{eq:markov-limit}
  \int_0^{t}\!\dd t_1\;(\cdots)
  \;\longrightarrow\;
  \int_0^{\infty}\!\dd\tau\;(\cdots)\,,
  \qquad \tau = t - t_1\,.
\end{equation}

\begin{keyresult}[Born--Markov master equation]
  Combining Assumptions~1--5, the reduced system state satisfies
  \begin{equation}\label{eq:born-markov}
    \eqbox{
    \frac{\dd}{\dd t}\tilde{\rho}_S(t)
    = -\int_0^{\infty}\!\dd\tau\;
      \tr_E\!\Bigl[\tilde{V}(t),\,
        \bigl[\tilde{V}(t-\tau),\,
              \tilde{\rho}_S(t)\tp\rho_E\bigr]\Bigr]}\,.
  \end{equation}
\end{keyresult}

\begin{remark}
  The intermediate equation before Assumption~5 (with
  $\tilde{\rho}_S(t)$ but the original limits) is the
  \textbf{Redfield equation}---time-local but with time-dependent
  coefficients.
\end{remark}

% ──────────────────────────────────────────────────────────────
\subsection{Example: two-level atom in the vacuum field}%
\label{sec:two-level-atom}

\begin{historical}[Spontaneous emission]
  Spontaneous decay of an excited atom was explained by
  Dirac--Wigner--Weisskopf shortly after the formulation of quantum
  mechanics---one of the earliest triumphs of quantum field theory.
\end{historical}

\subsubsection{The model}

The environment is a free bosonic field,
$H_E = \sum_k \omega_k\, b_k^\adj b_k$.
The system is a two-level atom with ground state $\ket{g}$, excited
state $\ket{e}$, and transition frequency $\omega_A$:
\begin{equation}\label{eq:H-atom}
  H_S = \omega_A\,\ketbra{e}{e}\,.
\end{equation}
The dipole coupling is
\begin{equation}\label{eq:V-dipole}
  V = \sum_k \bigl(g_k\,\sigma_+\tp b_k
      + g_k^*\,\sigma_-\tp b_k^\adj\bigr)\,,
\end{equation}
where $\sigma_+=\ketbra{e}{g}$, $\sigma_-=\sigma_+^\adj$, and the
coupling constants $g_k$ encode dipole matrix elements and mode
structure.

\subsubsection{Rotating-wave approximation}

In the interaction picture the coupling acquires phases
$e^{\pm i(\omega_A-\omega_k)t}$ (slowly rotating) and
$e^{\pm i(\omega_A+\omega_k)t}$ (fast counter-rotating).  Since
$\omega_A\sim 10^{15}\,\mathrm{s}^{-1}$, the fast terms
self-average and can be dropped (\textbf{rotating-wave
approximation}, RWA):
\begin{equation}\label{eq:V-int-atom}
  \tilde{V}(t)
  \approx \sum_k \bigl(
    g_k\,\sigma_+ b_k\, e^{i(\omega_A-\omega_k)t}
  + g_k^*\,\sigma_- b_k^\adj\, e^{-i(\omega_A-\omega_k)t}\bigr)\,.
\end{equation}

\subsubsection{Preview of the master equation}

Setting $\rho_E=\ketbra{0}{0}$ (vacuum) and carrying out the
Born--Markov procedure (detailed next lecture), one obtains:
\begin{equation}\label{eq:lindblad-atom-preview}
  \eqbox{
  \frac{\dd}{\dd t}\rho_S
  = -i[H_S,\,\rho_S]
  + \gamma\Bigl(
      \sigma_-\rho_S\,\sigma_+
    - \tfrac{1}{2}\sigma_+\sigma_-\rho_S
    - \tfrac{1}{2}\rho_S\,\sigma_+\sigma_-\Bigr)}\,,
\end{equation}
where $\gamma>0$ is the spontaneous-emission rate.

\begin{intuition}[Irreversibility from unitarity]
  The joint dynamics is unitary, yet the atom alone decays
  irreversibly.  The resolution: one excitation can flow into a
  continuum of field modes, but the reverse---all modes conspiring to
  re-excite the atom---has vanishing probability.
\end{intuition}

\begin{exercise}\label{ex:lindblad-form}
  Verify that~\eqref{eq:lindblad-atom-preview} preserves the trace
  and positivity of~$\rho_S$:
  \begin{enumerate}[(a)]
    \item Show $\frac{\dd}{\dd t}\tr(\rho_S)=0$.
    \item Writing $\rho_S$ in the $\{\ket{e},\ket{g}\}$ basis, derive
      the equations of motion for the matrix elements and check that
      $\rho_{ee}(t)=\rho_{ee}(0)\,e^{-\gamma t}$.
  \end{enumerate}
\end{exercise}
